% This is text associated with the open textbook "Open Electromagnetics"
% See immediately below for author, date
% This work is licensed under the Creative Commons Attribution-ShareAlike 4.0 International License. (CC BY-SA 4.0) 
% To view a copy of the license, visit http://creativecommons.org/licenses/by-sa/4.0/

%ORB TITLE Power Dissipated by an Electrically-Short Dipole
%ORB AUTHOR 0        # S.W. Ellingson, Virginia Tech (ellingson@vt.edu)
%ORB DATE 20191126   # yyyymmdd
%ORB ABSTRACT 

\label{m0208_Total_Power_Dissipated_by_an_Electrically-Short_Dipole} % <-- Must match name of this file and module directory name.  Don't mess with this!
\index{antenna!electrically-short dipole (ESD)|(}

The power delivered to an antenna by a source connected to the terminals is nominally radiated.
However, it is possible that some fraction of the power delivered by the source will be dissipated within the antenna. 
In this section, we consider the dissipation due to the finite conductivity of materials comprising the antenna.
Specifically, we determine the total power dissipated by an electrically-short dipole (ESD) antenna in response to a sinusoidally-varying current applied to the antenna terminals.  
(The ESD is introduced in Section~\ref{m0198_Radiation_from_an_Electrically-Short_Dipole}, and a review of that section is suggested before attempting this section.)
The result allows us to determine \emph{radiation efficiency} 
\index{radiation efficiency}
and is a step toward determining the impedance of the ESD.

Consider an ESD centered at the origin and aligned along the $z$ axis.
In Section~\ref{m0198_Radiation_from_an_Electrically-Short_Dipole}, it is shown that the current distribution is:
%
\begin{equation}
\widetilde{I}(z) \approx I_0 \left(1-\frac{2}{L}\left|z\right|\right) 
\label{m0208_eI}
\end{equation}
%
where $I_0$ (SI base units of A) is a complex-valued constant indicating the maximum current magnitude and phase, 
and $L$ is the length of the ESD.
This current distribution may be interpreted as a set of discrete very-short segments of constant-magnitude current (sometimes referred to as ``Hertzian dipoles'').
\index{Hertzian dipole}
In this interpretation, the $n^{th}$ current segment is located at $z=z_n$ and has magnitude $\widetilde{I}(z_n)$.

Now let us assume that the material comprising the ESD is a homogeneous good conductor.
\index{conductor!good}
Then each segment exhibits the same resistance $R_{seg}$.
Subsequently, the power dissipated in the $n^{th}$ segment is 
%
\begin{equation}
P_{seg}(z_n) = \frac{1}{2}\left|\widetilde{I}(z_n)\right|^2 R_{seg}
\label{m0208_ePseg1}
\end{equation}

The segment resistance can be determined as follows.
Let us assume that the wire comprising the ESD has circular cross-section of radius $a$.
Since the wire is a good conductor, the segment resistance may be calculated using Equation~\ref{m0159_eACZ2} (Section~\ref{m0159_Impedance_of_a_Wire}, ``Impedance of a Wire''):
%
\begin{equation}
R_{seg} \approx \frac{1}{2} \sqrt{\frac{\mu f}{\pi \sigma}} \cdot \frac{\Delta l}{a}
\end{equation}
%
where $\mu$ is permeability, $f$ is frequency, $\sigma$ is conductivity, and $\Delta l$ is the length of the segment.
Substitution into Equation~\ref{m0208_ePseg1} yields:
%
\begin{equation}
P_{seg}(z_n) \approx \frac{1}{4a} \sqrt{\frac{\mu f}{\pi \sigma}} \left|\widetilde{I}(z_n)\right|^2 \Delta l
\end{equation}

Now the total power dissipated in the antenna, $P_{loss}$, can be expressed as the sum of the power dissipated in each segment:
%
\begin{equation}
P_{loss} \approx \sum_{n=1}^N { \left[ \frac{1}{4a} \sqrt{\frac{\mu f}{\pi \sigma}} \left|\widetilde{I}(z_n)\right|^2 \Delta l \right] }
\end{equation}
%
where $N$ is the number of segments. This may be rewritten as follows:
%
\begin{equation}
P_{loss} \approx \frac{1}{4a} \sqrt{\frac{\mu f}{\pi \sigma}} \sum_{n=1}^N { \left|\widetilde{I}(z_n)\right|^2 \Delta l } 
\end{equation}
%
Reducing $\Delta l$ to the differential length $dz'$, we may write this in the following integral form:
%
\begin{equation}
P_{loss} \approx \frac{1}{4a} \sqrt{\frac{\mu f}{\pi \sigma}} \int_{z'=-L/2}^{+L/2} { \left|\widetilde{I}(z')\right|^2 dz' } 
\end{equation}
%
It is worth noting that the above expression applies to \emph{any} straight wire antenna of length $L$.  
For the ESD specifically, the current distribution is given by Equation~\ref{m0208_eI}.
Making the substitution:
%
\begin{flalign}
P_{loss} &\approx \frac{1}{4a} \sqrt{\frac{\mu f}{\pi \sigma}} \int_{z'=-L/2}^{+L/2} { \left|I_0 \left(1-\frac{2}{L}\left|z'\right|\right) \right|^2 dz' } \nonumber \\
         &\approx \frac{1}{4a} \sqrt{\frac{\mu f}{\pi \sigma}} \left|I_0\right|^2 \int_{z'=-L/2}^{+L/2} { \left|1-\frac{2}{L}\left|z'\right| \right|^2 dz' }
\end{flalign}
%
The integral is straightforward to solve, albeit a bit tedious.
The integral is found to be equal to $L/3$, so we obtain:
%
\begin{flalign}
P_{loss} &\approx \frac{1}{4a} \sqrt{\frac{\mu f}{\pi \sigma}} \left|I_0\right|^2 \cdot \frac{L}{3} \nonumber \\
         &\approx \frac{L}{12a} \sqrt{\frac{\mu f}{\pi \sigma}} \left|I_0\right|^2 \label{m0208_Ploss2}
\end{flalign}

Now let us return to interpretation of the ESD as a circuit component.
We have found that applying the current $I_0$ to the terminals results in the dissipated power indicated in Equation~\ref{m0208_Ploss2}.
A current source driving the ESD does not perceive the current distribution of the ESD nor does it perceive the varying power over the length of the ESD.
Instead, the current source perceives only a net resistance $R_{loss}$ such that
%
\begin{equation}
P_{loss} = \frac{1}{2}\left|I_0\right|^2 R_{loss}
\label{m0208_Pnet}
\end{equation} 
%
Comparing Equations~\ref{m0208_Ploss2} and \ref{m0208_Pnet}, we find
%
\begin{flalign}
\boxed{
R_{loss} \approx \frac{L}{6a} \sqrt{\frac{\mu f}{\pi \sigma}}  
}
\label{m0208_eRloss}
\end{flalign}
%
%%%%%%%%%%%%%%%%%%%%%%%%%%%%%%%%%%%%%%%%%%%%%%%
%ORB HIGHLIGHT
\begin{mdframed}[backgroundcolor=yellow!20,nobreak=true] 
The power dissipated within an ESD in response to a sinusoidal current $I_0$ applied at the terminals is $\frac{1}{2}\left|I_0\right|^2 R_{loss}$ where $R_{loss}$ (Equation~\ref{m0208_eRloss}) is the resistance perceived by a source applied to the ESD's terminals. 
\end{mdframed}
%%%%%%%%%%%%%%%%%%%%%%%%%%%%%%%%%%%%%%%%%%%%%%%
%
Note that $R_{loss}$ is \emph{not} the impedance of the ESD.
$R_{loss}$ is merely the contribution of internal loss to the impedance of the ESD.
The impedance of the ESD must also account for contributions from radiation resistance (an additional real-valued contribution) and energy storage (perceived as a reactance).
These quantities are addressed in other sections of this book. 

%%%%%%%%%%%%%%%%%%%%%%%%%%%%%%%%%%%%%%%%%%%%%%%
%ORB EXAMPLE
~\\ \begin{example} 
\begin{mdframed}[backgroundcolor=brown!20,linecolor=brown!20]\raggedright
Power dissipated within an ESD. 
\\~\\
A dipole is 10~cm in length, 1~mm in radius, and is surrounded by free space.
The antenna is comprised of aluminum having conductivity $\approx 3.7 \times 10^7$~S/m and $\mu\approx\mu_0$.
A sinusoidal current having frequency 30~MHz and peak magnitude 100~mA is applied to the antenna terminals.
What is the power dissipated within this antenna?

{\bf Solution.}
The wavelength $\lambda=c/f \cong 10$~m, so $L = 10~\mbox{cm} \cong 0.01\lambda$.
This certainly qualifies as electrically-short, so we may use Equation~\ref{m0208_eRloss}.
In the present problem, $a=1$~mm and $\sigma\approx 3.7 \times 10^7$~S/m.
Thus, we find that the loss resistance $R_{loss} \approx 9.49~\mbox{m}\Omega$.
Subsequently, the power dissipated within this antenna is 
%
\begin{equation}
P_{loss} = \frac{1}{2}\left|I_0\right|^2 R_{loss} \approx \underline{47.5~\mu\mbox{W}}
\end{equation}

\end{mdframed}
% note figures will need to go between "\end{mdframed}" and "\end{example}"
\end{example}
%%%%%%%%%%%%%%%%%%%%%%%%%%%%%%%%%%%%%%%%%%%%%%%

We conclude this section with one additional caveat:
Whereas this section focuses on the limited conductivity of wire,
other physical mechanisms may contribute to the loss resistance of the ESD.
In particular, materials used to coat the antenna or to provide mechanical support near the terminals may potentially absorb and dissipate power that might otherwise be radiated. 

\index{antenna!electrically-short dipole (ESD)|)}

